\clearpage
\SourceLeaf{241}%
\begin{center}
{\Large\textsc{[Bolyai and Lobachevsky.]}}



\emph{[\textup{\textsc{Gerling}} to \textup{\textsc{Gauss}}. Marburg, 18--21 December 1843.]}
\end{center}


\{\mbox{.\ .\ .\ .\ .\ .\ .} Two Hungarian scholarship holders from Maros Vásárhely brought a
letter of introduction from their university. It caught my attention because among
the other signatures I found the name \textsc{Wolfgangus Bolyai}. In the
past few days they have also handed me a private letter from another professor there
who had once attended my lectures. I took the opportunity
to ask about W. B. so that I could tell you about him. They describe him as a
rather stout, vigorous, and friendly old man who seems to be spending
a cheerful old age amid his gardens. His activity as a university
teacher is uninterrupted, hindered only by the fact that most listeners find it too
difficult to follow him. His son (of whom you told me that he had written a
good book) was pensioned off as a cavalry captain when he was an officer of only 32,
and lives with his father, if I understood correctly. --- Presumably
he thus gave offence with mathematical heresies. --- That former
student of mine promised at the time to arrange for me to obtain the books of
father and son here (for I once saw a book of the father's at your house too, which
you called interesting). He seems graciously to have forgotten it,
for his letter contains not a word about it. Nor can I find out the titles anywhere
(even in Berlin). --- Since the young people who have just arrived
are offering to attend to the matter, I should like to ask you
to have these titles written down for me when you have the opportunity. \mbox{.\ .\ .\ .\ .\ .\ .}\}


